% A 7-point stencil in 3D: computing one new cell reads its six face
% neighbours. Because the pattern is fixed and known, the cache behaviour can
% be predicted analytically rather than measured.
\documentclass[dvisvgm,border=3pt]{standalone}
\input{preamble.tex}
\begin{document}
\begin{tikzpicture}[x=1cm,y=1cm,
    cellb/.style={draw=uibkorange!70!black,line width=0.6pt,fill=uibkorange,
                  text=inkdark,font=\sffamily\scriptsize,
                  minimum size=7.5mm,inner sep=0pt},
    arr/.style  ={->,draw=blu,line width=1.0pt}]

  \def\d{1.15}

  \node[cellb] (c)  at (0,0)            {C};
  \node[cellb] (l)  at (-\d,0)          {L};
  \node[cellb] (r)  at (\d,0)           {R};
  \node[cellb] (t)  at (0,\d)           {T};
  \node[cellb] (bo) at (0,-\d)          {Bo};
  \node[cellb] (f)  at (-0.72,-0.72)    {F};
  \node[cellb] (ba) at (0.72,0.72)      {Ba};

  \foreach \n in {l,r,t,bo,f,ba}{\draw[arr] (\n) -- (c);}

  % the pattern is fixed, so the cache traffic can be modelled up front
  \draw[-{Stealth[length=5mm,width=5mm]},draw=blu,line width=2.6pt]
    (1.45,-0.95) -- (2.30,-1.55);
  \node[draw=uibkorange!70!black,line width=0.6pt,fill=uibkorange,text=inkdark,
        font=\sffamily\scriptsize,align=left,inner sep=5pt,anchor=north west]
    at (2.40,-1.35) {L1 misses: \#\\L2 misses: \#\\L3 misses: \#\\\ldots};

\end{tikzpicture}
\end{document}
